\section{物理起源}
\label{sec:physical-origin}

\subsection{量子场论}
\label{sec:QFT}

本节对（欧氏）量子场论（quantum field theory，QFT）作一个极其简短的介绍。

通常，一个量子场论具有如下输入数据：
\begin{enumerate}
	\item（\textbf{时空}）一个 $d$ 维黎曼流形 $(M,g)$；
	\item（\textbf{场}）$M$ 上的一个纤维丛 $B$，以及 $B$ 在 $M$ 上的截面空间 $\mathcal F=\mathcal{F}(M)=\Gamma(M,B)$；
	\item（\textbf{作用量泛函}）一个泛函 $S \colon\mathcal F \to \R$；
\end{enumerate}
非常粗略地说，应把 $\mathcal F$ 看成某个物理系统一切可能状态构成的空间，而函数 $S$ 则控制哪些状态更可能实际出现。

在经典物理系统中，我们想要考虑对物理量的测量，例如粒子的速度或位置。这可以用\textbf{可观测量}（observable）这一概念来形式化：按定义，可观测量就是一个泛函 $O \colon \mathcal{F}\to \R$。

其中特别重要的一类，是 $x$ 处的\textbf{局部算子}（local operator），它只依赖于场在 $x$ 处的取值或场在 $x$ 处的导数。

\begin{ex}[自由标量场论]\hfill 
\begin{enumerate}
	\item 一个（紧的）黎曼流形 $(M,g)$；
	\item $B=M\times \R$，从而 $\mathcal{F}=C^\infty(M)$；
	\item 由 $S(\phi)=\int_M  \phi\Delta_g \phi \mathrm{Vol}_g$ 给出的泛函 $S \colon\mathcal C^\infty(M) \to \R$，其中 $\Delta_g$ 是度量 $g$ 的 Laplace 算子，$\mathrm{Vol}_g$ 是 $g$ 对应的体积形式。
\end{enumerate}
当 $M=\R$ 时，对任意点 $x\in \R$，由 $O_x(\phi)= \phi(x)$ 与 $O_x^{(1)}(\phi) =  \phi'(x) $ 定义的泛函 $O_x, O_x^{(1)} \colon C^\infty(\R) \to \R$ 都是 $x$ 处的局部算子。

\end{ex}

另有两类场论对我们十分重要：
	\begin{enumerate}
		\item 设 $G_c$ 是紧 Lie 群。当 $\mathcal{F}$ 由 $M$ 上主 $G_c$-丛的联络组成时，这样的场论称为\textbf{规范理论}（gauge theory），$G_c$ 称为该理论的\textbf{规范群}（gauge group）。
		\item 设 $X$ 是流形。当 $\mathcal{F}$ 由 $M$ 到 $X$ 的映射组成时，这样的理论称为 \textbf{$\sigma$ 模型}（sigma model），$X$ 称为该 $\sigma$ 模型的\textbf{靶}（target）。此时 $B=M\times X$。  
	\end{enumerate}

物理学量子革命的一个洞见是：一个物理系统不能由单个场来描述，因为后者会对每个可观测量都取某个确定的值。我们只能把可观测量的期望值表示为积分，其中依赖于作用量的测度刻画了各状态出现的概率。这些积分常常形式地写作
\begin{equation}
    \langle O\rangle :=  \int_{\phi\in \mathcal{F}(M)} e^{ - S(\phi)/ \hbar  } O(\phi) D \phi.  \label{eq:expected-value}\notag
\end{equation} 然而，在许多情形下这些积分并无意义，因为空间 $\mathcal{F}(M)$ 往往是无穷维的，从而 Lebesgue 测度 $D\phi$ 无法定义。

另一方面，注意对 $M$ 的任意开子集 $U$，都可以定义其上的场空间 $\mathcal{F}(U)=\Gamma(U,B)$，从而也可以定义可观测量 $\obs(U) = \operatorname{Fun}(\mathcal{F}(U))$。\footnote{对任意 $U\subset V$ 都有 $\obs(U)\to \obs(V)$，于是 $x$ 处的局部算子空间 $\obs(x)$ 可等同于沿 $x$ 的全部开邻域 $U$ 所取的反向极限 $\displaystyle\varprojlim_{x \in U} \obs(U)$。} 在量子理论中，观测本身会扰动系统，因此人们无法在时空 $M$ 的同一点以相干的方式作两次观测。然而，我们仍有办法把不同点的观测结合起来：对 $V$ 内两个不相交的开集 $U_1$ 与 $U_2$，我们期待有一个映射 $\obs(U_1) \otimes \obs (U_2) \to \obs (V)$，它刻画\textbf{算子乘积}（operator product）：
\begin{equation*}
\tikz{
	\node[draw, circle, inner sep=0.5pt ] (a) at (-3.9,0) {$U_1$};
	\node[draw, circle, inner sep=2pt] (b) at (-3.2,0) {$U_2$};	
	\node[draw, circle, inner sep=15pt] (c) at (-3.5,0) {};	
	\node at (-3.5,0.5) {$V$}
}
\end{equation*}
更一般地，给定可观测量 $O_i$，它们只依赖于场在互不重叠的开集上的取值，我们考虑如下形式的积分
\[ \langle O_1,\cdots,O_n\rangle :=  \int_{\phi\in \mathcal{F}(M)} e^{ - S(\phi)/ \hbar  } O_1(\phi)\cdots O_n(\phi) D \phi.   \]	
它们称为该理论的\textbf{关联函数}（correlation function），是量子场论中的主要研究对象。

\subsection{拓扑量子场论}
\label{sec:TQFT}

在 Atiyah 与 Segal 的框架下，一个 $d$ 维\textbf{拓扑量子场论}（topological quantum field theory，TQFT）\footnote{对物理学家而言，TQFT 通常指关联函数在时空 $M$ 的度量 $g$ 连续变动时保持不变的量子场论。为区分这两个相关但不同的概念，我们把 Atiyah 与 Segal 意义下的 TQFT 称为「函子式 TQFT」（functorial TQFT）。在观念层面，只要有一个物理学家意义下的 TQFT，人们便可设想有一个函子式 TQFT，尽管实际构造一个可能相当艰巨。} 是从范畴 $(\operatorname{Bord}_d,\amalg,\emptyset)$ 到复向量空间范畴 $(\operatorname{Vect}_{\mathbb{C}},\otimes,\C)$ 的对称幺半函子 $Z$。$\operatorname{Bord}_d$ 的对象是闭定向 $(d-1)$ 维流形 $N$；从 $N$ 到 $N'$ 的态射是连接 $N$ 与 $N'$ 的 $d$ 维配边（bordism）$M$ 的微分同胚类；幺半结构由不交并 $\amalg$ 给出，单位对象是空集 $\emptyset$。

\begin{yizhu}[函子式 TQFT 的两条默认约定]
原文把 $\operatorname{Bord}_d$ 的态射取作配边的「微分同胚类」；严格说应取相对边界固定的微分同胚类，否则粘合不满足结合律。此外「幺半函子」意味着 $Z(N_1\amalg N_2)\cong Z(N_1)\otimes Z(N_2)$ 且 $Z(\emptyset)=\C$——原文随后把从 $\emptyset$ 到 $\emptyset$ 的配边 $M$ 的 $Z(M)$ 直接写成一个复数，用的正是 $Z(\emptyset)=\C$。
\end{yizhu}

把闭 $d$ 维流形 $M$ 视为从 $\emptyset$ 到 $\emptyset$ 的配边，便得到线性映射 $Z(M) \colon \C  \to \C $，亦即一个复数 $Z(M)$。物理上应设想 $Z(M)=\int_{\phi \in \mathcal{F}(M)} e^{ - S(\phi)/ \hbar  } D \phi $。另一方面，闭 $(d-1)$ 维流形 $N$ 所对应的复向量空间 $Z(N)$，正是该 TQFT 所描述的物理系统在 $N$ 上的状态 Hilbert 空间。

例如，先考虑 $d=1$ 的情形：我们得到一个称为「拓扑量子力学」的系统，其中 $H=Z(\operatorname{pt})$ 是系统单个粒子的状态空间，而由于理论的拓扑性质，时间演化是平凡的。在态–算子对应（state-operator correspondence）下，有 $Z(S^0)\cong \operatorname{End}(H)$，借助作用 $(a\otimes b)v=\langle b,v\rangle a$ 可把它与 $H\hat{\otimes}H$ 等同；物理学家更愿意把 $a\otimes b$ 写成 $|a\rangle\langle b|$，以提示这一乘法。

设 $d=2$。由于任何闭定向一维流形都是若干圆周的不交并，只需描述 $Z(S^1)$。此外，对应于「裤子」（pair of pants）的映射给出线性映射 $m\colon Z(S^1)\otimes Z(S^1)\to Z(S^1)$，对应于圆盘的映射则是线性映射 $u \colon  Z(\emptyset)\to Z(S^1)$：
\begin{center}
\begin{tikzpicture}[tqft/flow=east]
\begin{scope}[scale=0.4,every tqft/.style={transform shape},tqft/boundary lower style={draw}]
\node[tqft/reverse pair of pants, draw]   (a) {};
\path (a.outgoing boundary 1) +(-1,-2) node[font=\small] {\(\text{裤子}\)};
\end{scope}
\end{tikzpicture}\hspace{3em} 
\begin{tikzpicture}[tqft/flow=east]
\begin{scope}[scale=0.4,every tqft/.style={transform shape},tqft/boundary lower style={draw}]
\node[tqft/cap, draw]   (a) {};
\path (a.outgoing boundary 1) +(-0.4,-2) node[font=\small] {\(\text{圆盘}\)};
\end{scope}
\end{tikzpicture} \vspace{-0.5em}
\end{center}
事实上，下面两幅图说明 $m$ 与 $u$ 在 $Z(S^1)$ 上定义了带单位的结合代数结构：
\begin{center}
\begin{tikzpicture}[tqft/flow=east]
\begin{scope}[scale=1/4,every tqft/.style={transform shape},tqft/boundary lower style={draw}]
 \node[tqft/reverse pair of pants, draw, anchor=outgoing boundary 1] (a) {};
 \node[tqft/reverse pair of pants, draw, anchor=outgoing boundary 1] at (a.incoming boundary 1) (b) {};
 \node[tqft/cylinder to prior, draw, anchor=outgoing boundary 1]  at (a.incoming boundary 2) (c) {};
\path (a.outgoing boundary 1) +(1.5,0) node {\(=\)};
\path (a.outgoing boundary 1) +(1.5,-3.5) node[font=\small] {\(\text{结合律}\)};
\path (a.outgoing boundary 1) +(3,0) node[tqft/reverse pair of pants,draw, anchor=incoming boundary 1] (d) {};
 \node[tqft/reverse pair of pants, draw, anchor=incoming boundary 2]  at (d.outgoing boundary 1) (e) {};
 \node[tqft/cylinder to next, draw, anchor=outgoing boundary 1]  at (e.incoming boundary 1) (f) {};
\end{scope}
\end{tikzpicture}
\hspace{2em}  \begin{tikzpicture}[tqft/flow=east]
\begin{scope}[scale=1/4,every tqft/.style={transform shape},tqft/boundary lower style={draw}]
 \node[tqft/reverse pair of pants, draw, anchor=outgoing boundary 1]  (a) {};
 \node[tqft/cap, draw, anchor=outgoing boundary 1] at (a.incoming boundary 1)  (b) {};
 \node[tqft/cylinder, draw, anchor=outgoing boundary 1]  at (a.incoming boundary 2) (c) {};
\path (a.outgoing boundary 1) +(1.5,0) node  {\(=\)};
\path (a.outgoing boundary 1) +(4,-3.5) node[font=\small] {\(\text{单位公理}\)};
\path (a.outgoing boundary 1) +(3,0) node[tqft/cylinder,draw, anchor=incoming boundary 1] (d) {};
\path (a.outgoing boundary 1) +(6.5,0) node  {\(=\)};
 \path (d.outgoing boundary 1) +(7,0) node[tqft/reverse pair of pants, draw, anchor=outgoing boundary 1] (e) {};
 \node[tqft/cap, draw, anchor=outgoing boundary 1] at (e.incoming boundary 2)  (f) {};
 \node[tqft/cylinder, draw, anchor=outgoing boundary 1]  at (e.incoming boundary 1) (g) {};
\end{scope}
\end{tikzpicture} \vspace{-0.5em}
\end{center}
此外，乘法是交换的，因为可以交换乘法的次序。再把上面的图从右往左读，便可在代数 $Z(S^1)$ 上赋予交换 Frobenius 代数（commutative Frobenius algebra）的结构。事实上，$d=2$ 的 TQFT 由交换 Frobenius 代数分类。

注意，类似的想法适用于任何 $d$ 维 TQFT $Z$：利用从 $d$ 维球内部挖去两个不相交的 $d$ 维球所得的配边，可以证明当 $d\geq 2$ 时 $Z(S^{d-1})$ 也具有交换代数结构。

人们可能会问：从物理上讲，这个结构究竟从何而来？也就是说，我们固然知道 $Z(S^{d-1})$ 是 $S^{d-1}$ 上的状态空间，但一个 Hilbert 空间凭什么会带有乘法？

这一想法与量子场论中的局部算子概念有关。由于假设理论是拓扑的，$x$ 处的局部算子空间完全可以由 $d$ 维圆盘 $D^d$ 上的 $\obs(D^d)$ 忠实地表示。既然圆盘 $D^d$ 的边界是 $S^{d-1}$，就应把上面圆盘的图设想为给出一个线性映射 $\obs(D^d)\to Z(S^{d-1})$，其定义为 $O\mapsto O|0\rangle$：其中 $|0\rangle$ 是 Hilbert 空间 $Z(S^{d-1})$ 的真空态，局部算子 $O$ 被放在 $D^d$ 上的一点，而 $O|0\rangle$ 读出该算子所产生的状态。

在 TQFT 的语境下有一条非平凡的论断：线性映射 $\obs(D^d)\to Z(S^{d-1})$ 是同构，因此可以把 $Z(S^{d-1})$ 取作给定理论（在任意点 $x$ 处）的局部算子空间。这条论断称为\textbf{态–算子对应}（state-operator correspondence）。这就解开了谜团：$Z(S^{d-1})$ 上的交换代数结构，正是 TQFT 局部算子的算子乘积。同样清楚的是，「裤子」的图（以 $S^{d-1}$ 代替 $S^1$ 加以推广）也可以看作大圆盘 $D^d$ 内两个不相交圆盘 $D^d$ 的图，也就是上面算子乘积的图。

在物理学思想应用于数学的许多例子中，TQFT 的视角提供了一个有用的指导原则。在讨论如何运用这一思想之前，我们先解释如何从量子场论出发得到一个 TQFT。

\subsection{从量子场论到拓扑量子场论}

构造 TQFT（即不依赖于时空流形度量的理论）有两种众所周知的方式。其一是从不依赖度量的场空间与作用量泛函出发，例如 Chern–Simons 理论就是这样的理论。这一途径相当受限，能给出的例子相对较少。多得多的例子来自对（依赖度量的）超对称场论施加\textbf{拓扑扭转}（topological twist）。下面简要回顾后一想法。

考虑带标准度量的 $M=\R^d$。此时等距群 $\operatorname{ISO}(d,\R)=\operatorname{SO}(d,\R)\ltimes \R^d$ 称为 \textbf{Poincaré 群}（Poincaré group），它通过旋转与平移作用在 $\R^d$ 上。我们只考虑 $\R^d$ 上的场论，其作用量泛函在诱导的场空间作用下是等变的。

我们也只考虑截面空间 $\mathcal{F}$ 带有 $\Z/2\Z$-分次的场论；这在物理上源于粒子的自旋角动量，也正因如此，粒子自然地分为玻色子（偶）与费米子（奇）。

若一个场论容许非平凡的「奇对称」，就称之为\textbf{超对称}（supersymmetric，SUSY）场论，而这些奇对称称为\textbf{超荷}（supercharge）。

更确切地说，这意味着 $\mathcal{F}$ 带有一个 Lie 超代数 $\mathfrak{a}$ 的作用，该超代数称为\textbf{超 Poincaré 代数}（super-Poincaré algebra），其偶部是 Poincaré 代数 $\mathfrak{a}_0=\mathfrak{so}(d)\ltimes \R^d$，奇部 $\mathfrak{a}_1=\Sigma$ 由若干份 $\mathfrak{so}(d)$ 的自旋表示组成。Lie 括号由 $\mathfrak{so}(d)$ 在 $\Sigma$ 上的作用，以及 $\mathfrak{so}(d)$-表示的一个对称\footnote{在超 Lie 代数的世界里，若两个变元都是奇的，则 Lie 括号是对称的！}配对 $\Gamma \colon \Sigma\otimes \Sigma\to \R^d$ 给出。

为简单起见，从今往后我们都在超对称代数的复化上工作；就讨论扭转而言，这基本上是无害的。
\begin{ex}\hfill
\begin{enumerate}
	\item $d=2$ 的 $(\mathcal N_+,\mathcal N_-)$ 超对称代数 $\mathfrak{a}_{d=2}$，其奇部为 $ S_+\otimes W_+\oplus S_-\otimes W_-$，其中 $S_\pm \cong \C $ 是 $\mathfrak{so}(2)$ 的两个自旋表示，$\dim W_\pm= \mathcal N_\pm$。配对 $\Gamma \colon \Sigma\otimes \Sigma\to \C^2$ 由 $\mathfrak{so}(2)$-表示的同构 $S_+^{\otimes 2} \oplus S_-^{\otimes 2} \cong \C^2$ 诱导。
	\item $d=3$ 的 $\mathcal N=k$ 超对称代数 $\mathfrak{a}_{d=3}$，其奇部为 $S \otimes W$，其中 $S\cong \C^2$ 是 $\mathfrak{so}(3)$ 的自旋表示，$\dim W=\mathcal N$。配对 $\Gamma \colon \Sigma\otimes \Sigma\to \C^3$ 由 $\mathfrak{so}(3)$-表示的同构 $ \operatorname{Sym}^2 S\cong \C^3$ 诱导。
\end{enumerate}	
\end{ex}

最后，为了从超对称理论中提取 TQFT，设我们已选取超对称代数的一个超荷 $Q$，使得 $[Q,Q]=0$。由于 $Q$ 是奇的，这意味着 $\frac{1}{2}[Q,Q]=Q^2$ 在 $\mathfrak{a}$ 的任何表示上都作用为零。于是可以把 $\mathcal{F}$，甚至 $\mathfrak{a}$ 本身，视为 $\Z/2\Z$ 分次的复形，并取其 $Q$-上同调。这一手续必然得到一个更简单的理论，人们称之为\textbf{扭转}（twist，或称 twisting）。

若元素 $x\in \R^d\subset \mathfrak{a}$ 落在 $[Q,-]$ 的像中，则沿 $x$ 的平移在扭转后的理论中是平凡的。特别地，若 $O_i$ 是位于不同点 $x_i$ 的 $Q$-闭局部算子，那么只要这些点不相撞，沿 $[Q,-]$ 的像中的方向连续移动 $x_i$ 时，关联函数 $\langle O_1\dots O_n\rangle$ 保持不变。对我们最重要的情形是 $Q$ 的像填满整个 $\C^d$。此时对位置的依赖消失，理论变为拓扑的；因此，扭转后的理论称为原理论的\textbf{拓扑扭转}（topological twist）。

\begin{yizhu}[「扭转」兼指手续与结果]
「扭转」在本节兼指两件事：先指「取超荷 $Q$ 的上同调」这一手续，再指当 $Q$ 的像填满 $\C^d$ 时所得的拓扑理论。只有 $Q$-闭算子才在扭转后存活；沿 $[Q,-]$ 的像中的方向平移时，局部算子只改变一个 $Q$-恰当项，而后者对关联函数没有贡献，关联函数因而不变。
\end{yizhu}

拓扑扭转是否存在，完全取决于超 Poincaré 代数 $\mathfrak{a}$，从而取决于 $d$ 与 $\mathcal{N}$。关于这种对超荷数目的依赖，我们将用到两个重要事实。设 $\mathfrak{a}_{d=3}$（相应地 $\mathfrak{a}_{d=2}$）是 $d=3$（相应地 $d=2$）、$\mathcal{N}=n$（相应地 $\mathcal{N}=(n_+,n_-)$）的超对称代数。由标准的论证（例如见 \cite[\S\S 11.2 \& 12.1]{elliottTaxonomyTwists2022}），我们有：
\begin{itemize}
	\item 当 $d=3$（相应地 $d=2$）时，存在拓扑扭转当且仅当 $n\geq 4$（相应地 $n_\pm\geq 2$）。  
	\item 当 $n= 4$（相应地 $n_\pm= 2$）时，在适当的对称性意义下恰有 2 个拓扑扭转，我们将其记作 $Q_{\operatorname{A}}$ 与 $Q_{\operatorname{B}}$。
\end{itemize}

% [原文 p4–p8]
